\chapter{第七集 \quad 家国}

\section{正文}

陕西长武，泾水支流黑河的北岸上，碾子坡遗址笼罩在一片苍茫暮色中。这里，或许就是3000年前的豳，这里是黄土高原的南缘，千沟万壑分割出支离破碎的梁峁塬地，这里也是农耕区与半农半牧区的分界带。商以农为本，夷狄以牧为生，而周人的祖先曾摇摆于这两种生活方式之间，他们曾被戎狄侵扰，却又被商人当作戎狄一样征伐。商王武丁时期的甲骨上，历历可见与“伐周”相关的卜辞。古公亶父成为周族首领时，带着族人们从事农作，过上了家有余粮的安定生活。他常常站在豳原的高阜上南望，山势连绵，直接天际。他知道，翻过那一道道山梁，就是广袤的关中平原，那里，才是更适宜耕种的沃土，秋风萧瑟，熏育戎狄再次来犯，先掠才无，又意图占有土地和人口。

正当族人准备与戎狄决战时，古公亶父却制止了众人。他的理由是，对于普通人来说，在谁的治下生活并无不同，但战争却随时会令人丧生，因此他愿离开豳，将土地和人口都留给戎狄。

“古公亶父，来朝走马”，“率西水浒，至于岐下”。这是《诗经·大雅·緜》中周人对先祖自豳地迁往岐邑的记忆。古公亶父的退让，被豳人视为爱民的表现，纷纷离开家乡，扶老携幼，跟随他来到岐山之下。周边部族听闻后，也陆续前来归附，他们自此不再游牧，而在周原起城郭、筑屋舍、务农事，但这也意味着他们进入了“大邦商”的势力范围。商，实力强大，四处征伐，殷墟西北岗王陵中，小屯乙七建筑基址的祭祀坑，非王族的墓葬中，都发现了大量从葬的人骨，此时的周，不过是商邑西陲的蕞尔小邦，只能暂时的臣服于这个强大的王国。

与那个时代漠视生命的风气不同，姬姓周人的墓葬中，自始至终都极少有殉人，也几乎不见殉牲。多年来辗转求生的周族，深刻的认识到只有依靠群体的力量，才能存活于暑雨祁寒间，而仁德，就是凝聚起群体力量的黏合剂。古公亶父牵着小小的姬昌，在田间缓缓行走，耐心的教他辨识每一种农作物，娓娓讲述他们的习性，稚童眼神明亮，听得认真，努力记住祖父所说的一切，祖父常常讲起先祖后稷作农，并惠泽天下的故事。祖父还告诉他，人心，才是弱小的周族拥有的最大利器。夕阳一片，粟黍青青，祖孙俩脚下踏过的这片沃土，就是周原。周原遗址在今天的陕西扶风、岐山两县交界处，总面积约30多平方公里。诗经中说“周原膴膴，堇荼如饴”，意思是在周原这片丰饶的田野上，连生长出来的野菜也甘美如饴。

土地的膏腴不纯然来自天赐，也来自人力的开垦。从空中俯瞰，周原之上，周人曾以人工开挖的河道沟通了自然水系，以水库为调节枢纽，构成庞大的水利系统，灌溉着千里沃野。在周原的王家嘴，考古工作者们找到了灭商以前周人的世居之地。一号建筑基址正是那时的官室，它的总面积超过2200平方米，由前堂、后室、东西厢房、前后庭院和附属建筑组成。这种前堂后室的两进四合院式布局，明显以商的宫式建筑为模板。在王家嘴的东北方向，一水之隔处，发现了大小两重城垣，西北部的规整长方形小城，始建于商周之际至西周早期，大城城墙则是到西周晚期才扩建的，是万民所居的郭城。小城北部正中的凤雏建筑群可能是宫殿群，3号建筑基址是目前所知西周规模最大的单体建筑，在院落内发现了与土地祭祀有关的社石。

在这处建筑的南侧，还发现了一辆驷马驱驾的青铜轮牙马车，马匹及车辆装饰有大量镶嵌绿松石的青铜饰件，车厢外围装饰有薄壁青铜兽面、玉器、绿松石饰件。这辆华丽的西周马车，应该是高等级贵族的礼宾用车，以仁德和辑人心的周渐渐强大，在商王的授意下，周逐渐承担起为商治理西土的角色。姬昌的父亲季历，被商王帝乙封为西伯，率领周师四处征伐戎狄，并定期朝拜商王。仲春之月，渭水两岸桃花绽放，如云蒸霞蔚，河洲之上，雎鸠关关，众人簇拥着年轻的姬昌从奇异的宫殿中迤逦而出，来到渭水边，河对岸走来了他的新婚妻子太姒。周族传统的联姻对象是同样生活在商人西土的姜族，传说中正是姜族的女子姜嫄，诞育了周族的祖先后稷。姬昌的祖母也是姜族的太姜。

但自姬昌的父亲季历起，周族首领也开始与东方诸族联姻，试图融入中原社会。姬昌的母亲太任，是殷商的贵族挚任氏，姬昌的妻子太姒也从东来。凤雏建筑遗址出土了17000多片卜甲卜骨，其中有字卜甲有190多片，字体纤细如微雕，所书文字与商人文字无异，所记载的占卜程序也与商人相同。卜辞多属姬昌时期，内容包括向殷商王室的祖先成汤、太甲、文丁、帝乙等祭祀，其中还出现了周方伯的称呼，这些都足见周人对商的顺服。另一方面，姬昌继承了古公亶父以来的理念，笃仁、敬老、慈少、礼贤，前来归顺于周的人越来越多。岐山苍翠的林中，几声清啼打破清晨的寂静，一只锦雉掠过晨雾落在树梢上，而洹水之畔的大邑商，商王帝辛彻夜未眠，他摒弃了职业的贞人，亲自手持甲骨，细观纹理，推究上帝和祖先的启示。

突然，一阵鸟叫声打破了他的迷思，他惶惑地望向窗外，周的强大终于引起了商人的警惕，姬昌被帝辛囚禁于羑里七年，经周臣的多番奔走才终能得释。姬昌归去后，一面继续修德政，一面为商王征伐叛邦，赢得越来越多诸侯的归附，但很快姬昌在迁至丰京的次年死去。姬发即位后，修建了镐京。丰镐遗址在今西安市长安区，沣河纵贯两京之间，丰京在河西，镐京在河东，隔河连片，蔚然可观。镐京遗址内，已发现了十余座大型夯土建筑，还有道路和排水设施。郿坞岭高地14号建筑的规模宏伟，附近还有瘗埋动物骨骼的祭祀坑，或许就是周人的宗庙。周人以血缘确定贵族的亲疏、等级、分封和世袭关系。姬发居于宗周，也就是丰镐，他是姬姓宗族的族长，也是国君，掌握着族权和国家最高权力。

都邑之外，环绕着王族子弟的采邑。这里是国也是家，有国之架构，也有家之血缘。周都的东迁和规模扩大，无不宣示着这个曾经的西陲小邦如今的壮大。他们披仁德之坚，执宗法之锐，酝酿着对东边那个大邦的奋力一击。暮色渐深，马车仍在山路间向残阳落下的方向奔行，艰难的路途如同鸿沟，隔离了远方的山河。太师疵和少师彊两位商的重臣牢牢护着身旁的编铙，这是商王朝最神圣的祭乐之器，也是他们家族世代职责所司。他们生怕一路的颠簸令其稍有毁损。由于帝辛的残暴，他们不得不和许多殷商贵族一样，选择离开故土与故国，逃往以仁德闻名的周。在伐商的檄文《牧誓》中，武王痛责纣王的失德，他将克商之战称为“行天之罚”。

在那个极度崇拜神巫的时代，周人却从世俗层面，提出统治的合法性应该来自天命，而天命不在于祖先与天的特殊关系，不在于独占通天之技术和资源，主要在于统治者的德。有德之君，就会被天授予神圣的大命，得此大命者，可以革除失德王朝的旧命，建立新的王朝，抚育万民。陕西临潼零口镇出土的利簋，其底部铭文记录了武王伐商的事件，“珷征商，佳甲子朝，岁鼎，克昏夙有商”，意思就是甲子日的清晨，时逢岁星当空。武王的军队仅用一日便结束了战争。身为右使的利，由于跟随武王作战得胜后受到奖赏，因此铸造了这件青铜簋用以记功，并作为祭奠祖先的礼器。考古学家根据铭文中提及的特殊天象，再加上对干支纪年的推算，认为武王伐商的具体年代可能是在公元前1046年。

商，以灿烂的青铜和完善的文字体系，将我们的中国文明推送到新的发展高度。但是商的覆亡也说明，对神权的独占以及大规模的征战和杀殉，已经不足以威服四方。面对辽阔的国土和多元的族群，如何将之真正纳入到统一中央权力的管辖之下，周人需要新的政治方略。周人代商而立后，妥善安置了殷商遗民，帝辛之子武庚被封于殷商故都，仍可延续殷祀。但同时，武王又将自己的三位胞弟管叔、蔡叔、霍叔分封在殷都附近，起着监控的作用。他最信任的4位重臣则被分封在原商朝疆域的东、北、南境。姜尚在齐国，周公在鲁国，召公在燕国，南公在曾国，对商构成巨大威胁的东夷、淮夷曾活跃在这些地方，武王不敢掉以轻心。这一年，箕子从东边浮海而归，前往周都朝见周王。

归途中，他特地折往商人旧都，只见宫倾城圮的废墟上，已是禾黍青青，箕子是殷商王族，也是与微子比干齐名的重臣。曾经的朝堂之上，面对一意孤行的帝辛，他没有像比干一样死谏，也没有像微子一样离去，而选择了装疯，终致被囚为奴。武王之前，他也曾面陈《洪范》献上治国之法，但终不愿仕周为官而远走。眼前，太行西峙，洹水东流，故景依旧。那个曾令天下“莫敢不来享，莫敢不来王”的天邑商，却已无迹可寻。箕子忍不住悲从中来，在长天后土之间，怆然的放声唱道，“麦秀渐渐兮，禾黍油油，彼狡童兮不与我好兮”。诗中的狡童，正是商朝的末代之王，他的侄子帝辛，早已在武王兵临城下那日自焚于鹿台之上。殷商遗民听到这支歌，莫不垂泪。他们的王朝啊，终是沉入历史的长夜了。

月华如洗，周王的寝殿中依旧灯火通明，武王克商返回镐京后，却陷入比出征前更深的忧虑中，这天，他又再一次的熬到长夜欲曙，忍不住派人换来了周公。武王和周公的彻夜长谈，牵涉到周人疆域和身份观念的转型。周人本来自认为是西土之人，而商人所在则为东国，全部占有商王朝的原有领土后，周人的志向并不仅仅是取商人而代之，还希望能建立一个规模空前的王国。一番长谈后，二人商定，要在伊洛二水之北营建洛邑，以“定天保，依天室”，天室，就是太室山。太室山为天下之中的观念由来已久，夏商两代的都城都在这附近。武王返回镐京途中，曾在姜尚的陪同下，登上太室山山顶，俯瞰四方山川形势，那时他大概就已经存在营建洛邑的念头了。

这是何尊，出土于距周人旧都岐邑不远的陕西宝鸡贾村镇。尊，是酒器，同样也是礼器，这件青铜器是武王之子周成王在位第5年时，周的宗室贵族何，为纪念成王的训诫和封赏而作。何尊高38.8厘米，重达14.6公斤，内底铸有12行122字铭文，铭文中提及，成王将政治中心迁往成周时，追述了武王克商后的敬告上天之辞。“余其宅兹中国，自之乂民”，这是我们目前看到的最早的以文字书写的“中国”之名。成王即位后，在周公、召公的辅政下，完成了武王营建洛邑的遗愿，自此后，丰镐为宗周，洛邑为成周，也即是中国。在中国立都，周人便不再是西土之人，周也不再是西土方国，而是位居天下之中的宗主国。普天之下，莫非王土？这里是周人以天下视角确立的自然地理和政治地理核心。

周公旦在余生中仍常常想起他与哥哥姬发的那次彻夜长谈，除了营建洛邑之外，那时已知道自己病体难支的武王，还提出了另一个请求，若自己身死，便请周公继位。震惊之余，周公答道，周人的立族之本乃是宗法，族权自来是父死子继，国，亦当如此。王权的子继之法一旦实行，嫡庶之制、宗法之制、分封之制，便随即顺理成章，将密织成一整套坚固的上层建筑。周公所思，武王岂能不懂？身为国君，他惟愿自己死后仁德如周公者，能接过这未竟的事业，带领根基未稳的周，真正成为长久存在的大邦。周公的心术与规摹，并不在于眼前一时的权利，而在于周的万世治安大计。因此他选择了不登王位，不就封国，只留在王畿辅佐年少的成王，平武庚之乱东征淮夷，劳心竭力地辅政，敦促成王立德修身，既是为宗族，亦是为天下。

在周公的辅佐下，周王朝一面牢牢掌控着王畿之地，一面通过分封的诸侯国于四方之土，探索着不同的治理模式，使各地的不同部族，得以融入周人的政治秩序。鲁国的伯禽，面临着相当复杂的局面，当地遍布东方诸族，文化迥异。当看到熟悉的家门时，伯禽忍不住加快了步伐，走入堂中。他离开父母前往鲁地就封已经整整3年。因为父亲周公要留在成王身旁辅政，便只能由他代为就封。此番借着回来报政，他终于可以与父母相聚几日，堂上的周公白发暗生，面容因为常年的殚精竭力而略显疲惫，但双目依旧精光四射，仿若能洞悉世间的一切事理。他和蔼而专注的听伯禽汇报着三年来在鲁国的治理成果，频频颔首。末了，周公似乎想起了什么，问伯禽为何3年才回来报政？伯禽早有准备，从容答道，因为需要变其俗，革其礼。

周公告诉伯禽，太公姜尚就封鲁国之邻的齐国后，精简君臣之礼，依从当地习俗，仅过5月便回来报政。伯禽愣住了，周公告诫道，只有所施政策平易近民，才能得民心归附。周公所称许的齐国，是周初三公之一太公姜尚的封地。齐地虽有工、商、渔、盐之饶，却也是直接面对东夷诸部族的最前线。近年在山东临淄高青发现一处西周墓地，墓葬规格很高，也有祭祀遗迹和车马坑。大墓内出土的铜器，出现了齐公字样的铭文。若这里是齐公墓地，那么或许最早的齐国都城便在这一带。莒县西大庄春秋墓中出土的一件铜甗，可能是齐侯嫁女的陪嫁。莒是东夷势力之一，藉由此器，可以想见当时齐国通过各种手段融合东夷的努力。燕国所在之地，是周人克商后新辟的封地，为周王朝守卫着北境，当地同时还生活着迁徙而来的殷移民以及本地的土著族群。

这是北京房山琉璃河遗址出土的克盉，铭文“命克侯于燕”一句，意指周王将克封于燕，治理这里的疆土和六族人民。克，是太保召公奭的长子，因为召公留在王畿辅佐周王，故由长子代为就封。然而，1902号墓中出土的作册奂卣中太保墉燕的铭文，说明召公本人不仅亲自来过这里，还主持了燕都的营建。1902号墓的墓主人是一位殷遗民，他与周人同葬于一片墓地内。这位名叫奂的史官，生前颇受重用，随葬的青铜器上，铭刻着曾受周人封赏的荣光。他规矩的沿袭着商人殉狗的旧习俗，却也认同于周人建构的新秩序，由此可见。周人在族群面貌复杂之地采取的并存策略，收效应该比较显著。与燕地一样，汉水流域也是周王朝灭商后新扩的领土。

考古工作者在湖北随州附近，发现了一个失载于文献的曾国，曾之始祖南宫适也属姬姓，是辅佐文王和武王完成灭商大业的重臣，因为本人要留在王畿领职，只能由儿子代为就封。叶家山墓地中发现了曾侯犺、曾侯谏等数代西周曾侯的墓葬。这个从未出现在任何传世文献中的国家，现在人们普遍相信，他就是左传中汉阳诸姬之首的随国，一直存在了700余年，留下众多两周时期的珍贵文物和历史信息。这件方鼎是叔虞自作的铜器，叔虞是武王之子，成王之弟，亦在成王时期获封唐地，留下了桐叶封弟的传说。叔虞之子燮父继位后，迁都晋水之傍，将国号自唐改为晋。晋南是夏人故地，自然启以夏政，那里又遍布戎狄部族，亦须疆以戎索，融合各方，正是晋国肩负的重要责任。

燮父墓中出土了一件精美的鸟尊，器身造型生动无比，通体纹饰华美，堪称西周青铜艺术的杰作。鸟尊盖内铭文中的晋侯，就是燮父的自称，这个晋字见证着燮父迁唐于晋的历史。在山西曲沃曲村天马遗址的北赵晋侯墓地中，发现了9组19座排列规整的晋侯夫妇墓及附属的车马坑，视为国君夫妇的专属公墓。每代晋侯夫妇几乎都是一夫一妻异穴合葬，唯有64号墓中的晋穆侯，身旁并穴而葬了两位夫人。其中一位夫人墓内随葬的组玉佩，是迄今所见最复杂的。周代组玉佩是王室贵族们进行祭祀、朝聘、婚配等活动时佩戴的仪礼重器，不但规范了君臣尊卑，也将周人的道德观念赋予其上，是西周社会礼制日臻完善的重要体现。成年的姬诵头戴冠冕，身着朝服，肃容立于大朝会坛场上。周公在左，太公姜尚在右，皆衣冠济楚。

堂下为首的是尧舜禹商后裔，他们的出现象征着周的合法王统。诸侯们依次献上国内珍产，东方的海蛤、玄贝，南方的翠鸟、大象，西方的孔雀、朱砂，北方的麋鹿、大鲵，琳琅满目的异产，象征着各地对周的臣服。这场隆重的大朝会礼仪，从时间和空间两个维度上，宣示着周天子的权威。周初共分封71国，其中与周王族同姓的姬姓就有53国。分封，是周人总结商人治理得失，结合自身文化传统的创造。这项制度的核心，不是强力的军事控制，而是不同族群在周人统领下的融合，外封的诸侯都承担着以周文化融合当地族群的使命。周人将根植血缘关系的宗法制与国家治理紧密联系，使得家国一体。凭借稳定的社会体系，周天子的礼仪政令得以有效传递，周王朝的国家机器得以有效运行。

陕西扶风庄白一号窖藏，是新中国成立以来，出土青铜器数量最多的窖藏，在面积仅2平方米的窖穴内，就收纳了103件青铜器，其中74件都有铭文。在这件共王时铸造的史墙盘上，追述了文武之后迄至穆王的各代周王功绩。成王严明纲纪，封建邦国，康王开辟疆域，昭王广拓荆楚，穆王深谋大略。铭文在追述周人世系时，最大的变化，是直接以始祖后稷与上帝并列，称上帝后稷亢保。这意味着，周人将本族的祖先后稷，往前推至于舜禹同时，将族史与国史交织，逐渐融入正统华夏的叙事。这一古史系统的完成，应是在众多变革发生的穆王时期。夏将至，伯中受穆王之命，率兵在棫林追击淮夷，最终克敌制胜，穆王后王俎姜派人来到前线对他进行赏赐。凯旋回朝后，他便马上赶往故乡告祭父母。父亲在他幼年时便已辞世，是母亲将他抚育成人。

以往每次出征前，母亲总会端正他的衣冠，细细检查他的甲胄是否有破损，嘱咐他既要勇敢战斗，也要保全自己。淮夷之战的不久前，母亲刚刚去世，再无人那样絮絮地叮咛。此刻，他在心里默默告诉母亲，他平安归来了。伯中为母亲所做的簋上，铭刻着一个儿子对亡母的殷殷思念之情。“母亲的美好品行，在我心中涤荡，永远护佑我身，让我战无不胜，战毕，我身无伤”。同墓所出还有伯中为父母所做的两件方鼎，一件记述了穆王之命，一件记述了穆王后之赐，从周穆王、共王开始，伯中墓中这样鼎、簋配合的食器，成为了新的礼器体系主流商人创立的以觚、爵等酒器为主的青铜礼器系统被逐渐摒弃。周人认为，商人纵酒是致使其亡国的因由之一。

从武王的《牧誓》到周公的《酒诰》，从西周早期康王的大盂鼎，到西周末年宣王的毛公鼎的铭文中，都一再警示周人上下勿要酗酒，或许因此，周人才逐渐创立了自己的礼器体系。此时，青铜礼器的风格也发生了改变。商代晚期以来，华丽繁缛的青铜器装饰风格逐渐被简化，质朴简洁的鸟纹、几何纹逐渐成为主流。分封制构建的家国体系、以德配天的政治理念被寄寓在这些世俗色彩浓厚的青铜礼器中，物化成为礼制。周人之礼，如春雨润物无声，如阳光普照四方，衣食住行、婚丧嫁娶，家国朝野，无处不在。墓葬中的青铜礼器，便是这套礼制最直观的表现。根据《春秋公羊传》的记述，在礼仪活动中，天子使用九鼎。在他之下的诸侯、大夫、元士，则依次使用七鼎、五鼎、三鼎，这套礼器制度，在三门峡上村岭虢国墓地中展露无遗。

这些自西周末年延续至虢国灭亡的墓葬中，以食器、车马器、乐器的数量和有无分为五等。上世纪90年代，虢国墓地进行了大规模考古工作，发现了虢季墓和虢仲墓两座国君级大墓。虢仲墓随葬各类鼎20余件，其中包括一套大小成列的七鼎，与其王室卿士的身份相符。第二等级之下的墓葬皆无乐器，其余器类的数量逐级递减，五个等级之间，不止青铜礼器的种类和数量有别，食器所盛菜肴，车马所载出行，黄钟所奏雅乐更是不同。这百余年间，数百座墓葬都遵守如一的规则，正是雍雍穆穆的周礼。清道光末年出土于岐山的毛公鼎，是毛公为纪念周宣王的册封和赏赐而作。铭文中宣王先是追怀文王、武王时以仁德肇国，君臣一体治理天下的盛世，继而感怆时艰，为局势不宁、国家濒危而忧心忡忡。

由是恳切请求身兼重臣和宗亲的毛公，要为我邦我家而操心劳力，需勤理政事，勿要搪塞庶民，勿让官吏徇私，勿令鳏寡失养，望能约束僚属，勿使酗酒，你切莫失职，要夙夕谨记守业不易。从宣王的话语中，那个时代的风雨如晦扑面而来。西周王朝自始建后，在相当长的一段时间里，都将疆土经营的重心放在东土、北土和南土，而对宗周所在的西土，或许正因为是本族发祥之地，反而自信的认为能将其置于掌控下，这样的疏于防备留下的隐患。宣王时期的短暂中兴，终究无力回天。至幽王时，这种东西失衡的战略布局终于彻底暴露其弊端。申侯以女儿申后和外孙宜臼被废一事为由，联合缯国与犬戎，杀幽王于骊山下。此后，平王迁往成周，数百年来的祖居故地不得不被放弃。

在周原一带，发现了西周晚期的百余处青铜器窖藏，其中出土铜器数量较多，且有长篇铭文的重要发现就有十几处。当时这些铜器群应该分别被不同的贵族家族所拥有。然而他们为何几乎在同一时间段将这些重器藏于地下，最终又未再取出？有人联系西周末年的局势，推测可能是与犬戎入侵有关，贵族们仓促掩埋了难以带走的宝器，匆匆跟随周王东行。此后，由于各种原因，致使他们就此睽违天日，长埋地下数千年，成为后人了解一部煌煌西周史的锁钥。在周由盛转衰的同时，一支此前默默无闻的西陲小族，进入了周王室和东方诸侯的视线，那便是秦。秦人早期的历史，几乎重复着周人先祖的道路，从艰难求存于西戎之侧，到为周王朝守卫西土。幽王被杀时，秦襄公曾带兵相救，虽未能力挽狂澜，却得以崭露头角，后又以兵护送平王东徙成周，平王也因此正式锡封襄公为诸侯，并赐以岐西之地。

后来，春秋时秦武公所铸的秦公镈上便有追忆此事的铭文，“我先祖受天命，赏宅受国”。自此之后，西边的秦国开始崛起于汧渭之间，那时候，谁也不知道他们即将成为下一个时代的主角，并开启中国历史上一段浓墨重彩的篇章。孔子曰，郁郁乎文哉，吾从周。周人通过新的政治制度和政治理念的构建，形成了中华文明的人文基调，在随后的春秋战国时代焕发出更璀璨的光辉。华夏这一文明体的概念自此确立，并开启了一种内含天下结构的政治共同体形成模式。从此，不论是何种语言、文化或生活方式的族群，不论从何处发源，只要在孕育华夏文明的土地上生活，通过不断的调适认同基础，终将能融入这套无外的天下秩序里。那最初在伊洛之间的中国，势必如滚雪球般，生长成一个灿烂而盛大的政治共同体！
